\chapter{Laboratory 1: IPv6 Addressing and Static Routing on Cisco Routers}

\section{Objectives}

This laboratory builds a small IPv6 network made of five Cisco routers in GNS3 and configures it so that every router can reach every other router. Three mechanisms are covered: enabling IPv6 routing on a router, assigning IPv6 addresses (both manually and through stateless autoconfiguration), and configuring static routes. At the end, the connectivity is tested with \texttt{ping}, and the control messages used by the protocol (Router Solicitation, Router Advertisement, Neighbor Solicitation and Neighbor Advertisement) are observed with Wireshark.

The laboratory runs in the GNS3 virtual machine. The installation procedure is described in the file ``Lab intro - GNS3.pdf'', in the ``Material'' section of the course.

The new concepts introduced here are summarized in Section~\ref{sec:lab01-glossary}, and the commands are collected in Section~\ref{sec:lab01-commands}.

\section{Network topology}
\label{sec:lab01-topology}

The network is made of five routers, R1 to R5, connected as shown in Figure~\ref{fig:lab01-topology}. R1 is an edge router connected only to R2. R2 is the central router: from it two parallel paths start, one through R3 and one through R4, and both paths end at R5. The network therefore has two possible routes between R2 and R5. This redundancy is what makes the routing configuration interesting, because each router must know which of the two branches leads to each destination.

\begin{figure}[h]
\centering
\safeimg[width=0.85\linewidth]{media/lab01_topology.png}
\caption{Network topology used in the laboratory (Figure 1 of the handout).}
\label{fig:lab01-topology}
\end{figure}

The cabling must be reproduced exactly, because the addresses are tied to the physical links. Table~\ref{tab:lab01-cabling} lists the links.

\begin{table}[h]
\centering
\begin{tabular}{@{}llll@{}}
\toprule
\textbf{Link} & \textbf{Router} & \textbf{Interface} & \textbf{Connected to} \\
\midrule
R1 -- R2 & R1 & FastEthernet0/0 & R2 FastEthernet0/0 \\
R2 -- R3 & R2 & FastEthernet0/1 & R3 FastEthernet0/0 \\
R2 -- R4 & R2 & FastEthernet1/0 & R4 FastEthernet0/0 \\
R3 -- R5 & R3 & FastEthernet0/1 & R5 FastEthernet0/1 \\
R4 -- R5 & R4 & FastEthernet0/1 & R5 FastEthernet0/0 \\
\bottomrule
\end{tabular}
\caption{Physical links of the topology.}
\label{tab:lab01-cabling}
\end{table}

Notice that on R5 the order is reversed with respect to what one might expect: FastEthernet0/1 goes to R3 and FastEthernet0/0 goes to R4. This is visible in the GNS3 screenshot of the topology, and it is the wiring that the addresses below assume. Before assigning any address, it is worth checking the wiring in GNS3 once more. The addresses of each link are fixed by the naming convention and by the cabling, so an address placed on the wrong interface makes the routing fail even when every command is syntactically correct.

\section{Addressing plan}

The laboratory assumes that IANA has assigned the prefix \texttt{2001:1000::/56} to the network administrator. Designing the addressing plan means dividing this prefix into one subnet per link and assigning an address to every router interface.

\subsection{How the prefix is divided}

A \texttt{/56} prefix fixes the first 56 bits. The remaining 72 bits are the host part, but the protocol fixes the last 64 bits as the interface identifier, so the bits left for subnetting are bits 57 to 64. These 8 bits are the low byte of the fourth hextet of the address. The high byte of the same hextet is part of the \texttt{/56} prefix and is therefore zero. With 8 subnet bits there are $2^8 = 256$ possible \texttt{/64} subnets, which is more than enough for five links.

A \emph{hextet} is one of the eight 16-bit groups of an IPv6 address, written as four hexadecimal digits and separated by colons. In \texttt{2001:1000:0:12::/64} the four hextets written before the double colon are \texttt{2001}, \texttt{1000}, \texttt{0} and \texttt{12}; the double colon stands for the remaining zero hextets.

\subsection{Naming convention}

The laboratory uses a simple convention that makes the addresses easy to read: the subnet of a link is \texttt{2001:1000:0:xy::/64}, where \texttt{x} and \texttt{y} are the numbers of the two routers connected by the link. Each router, on each of its links, uses the host identifier \texttt{::n}, where \texttt{n} is its own router number. For example, R2 on the R1--R2 link has the address \texttt{2001:1000:0:12::2/64}.

The router numbers are 1 to 5, so writing them as hexadecimal digits is unambiguous: the digits are the same in decimal and in hexadecimal.

The \texttt{/64} length is not a free choice. Stateless autoconfiguration, which R1 uses, builds the address by appending a 64-bit interface identifier to a 64-bit prefix. For this reason \texttt{/64} is the standard length for every link with hosts that autoconfigure their addresses.

Table~\ref{tab:lab01-plan} shows the complete plan. R1 is not in the table with a manual address, because its address is generated automatically (Section~\ref{sec:lab01-autoconfig}).

\begin{table}[h]
\centering
\footnotesize
\begin{tabular}{@{}p{3.4cm}p{1.2cm}p{5.2cm}p{5.2cm}@{}}
\toprule
\textbf{Subnet} & \textbf{Link} & \textbf{Router A interface} & \textbf{Router B interface} \\
\midrule
\texttt{2001:1000:0:12::/64} & R1--R2 & R1 Fa0/0: autoconfigured & R2 Fa0/0: \texttt{2001:1000:0:12::2/64} \\
\texttt{2001:1000:0:23::/64} & R2--R3 & R2 Fa0/1: \texttt{2001:1000:0:23::2/64} & R3 Fa0/0: \texttt{2001:1000:0:23::3/64} \\
\texttt{2001:1000:0:24::/64} & R2--R4 & R2 Fa1/0: \texttt{2001:1000:0:24::2/64} & R4 Fa0/0: \texttt{2001:1000:0:24::4/64} \\
\texttt{2001:1000:0:35::/64} & R3--R5 & R3 Fa0/1: \texttt{2001:1000:0:35::3/64} & R5 Fa0/1: \texttt{2001:1000:0:35::5/64} \\
\texttt{2001:1000:0:45::/64} & R4--R5 & R4 Fa0/1: \texttt{2001:1000:0:45::4/64} & R5 Fa0/0: \texttt{2001:1000:0:45::5/64} \\
\bottomrule
\end{tabular}
\caption{Addressing plan of the laboratory.}
\label{tab:lab01-plan}
\end{table}

\section{Step 1: Resetting the routers}

The routers used in the laboratory may already contain configuration from earlier attempts. A leftover command, such as a wrong address or a route with a typo, can make the final result impossible to interpret, so the first action is to bring every router back to a clean state.

On a router that has unsaved changes, the simplest reset is a reload that discards them. In privileged mode (the prompt \texttt{R1\#}, reached with \texttt{enable}), the command \texttt{reload} restarts the router. IOS asks whether to save the current configuration; answering \texttt{no} discards the changes made since the last save. To erase also the saved configuration, the command \texttt{erase startup-config} is run before \texttt{reload}. The router then starts with no configuration at all.

\begin{codeblock}
enable
reload
\end{codeblock}

\begin{codeblock}
enable
erase startup-config
reload
\end{codeblock}

After the reset, the command \texttt{show running-config | include ipv6} must print nothing. The vertical bar is a filter: it shows only the lines of the configuration that contain the word \texttt{ipv6}. An empty output confirms that IPv6 is not configured yet.

\section{Step 2: Building the topology in GNS3}

In GNS3 five routers are added to the project and connected according to Table~\ref{tab:lab01-cabling}. Each link is a cable between two free interfaces, one on each of the two routers. The routers are then started, and the console of each one is opened to type the commands.

The interface labels shown in GNS3, such as \texttt{0/0}, correspond to the IOS names \texttt{FastEthernet0/0} used in the console. The router names must also match the figure, since the addresses and the routes of this laboratory refer to R1 to R5 by number.

\section{Step 3: Enabling IPv6 routing}

On each of the five routers, IPv6 must be enabled before anything else. The commands are the same for all routers:

\begin{codeblock}
enable
configure terminal
ipv6 unicast-routing
end
\end{codeblock}

The three modes of the IOS command line are visible in the prompt. \texttt{R1>} is user mode, where only read-only commands are available. \texttt{R1\#} is privileged mode, reached with \texttt{enable}, where the \texttt{show} and \texttt{reload} commands are available. \texttt{R1(config)\#} is global configuration mode, reached with \texttt{configure terminal}, where the commands that change the router are entered. The command \texttt{end} returns to privileged mode.

The command \texttt{ipv6 unicast-routing} has two effects. It enables the IPv6 protocol on the router, and it allows the router to forward IPv6 packets between its interfaces. Without it, a router can have IPv6 addresses but will not route IPv6 traffic. The handout requires this command to be the first IPv6 command typed on the router, before any address or route is configured.

\section{Step 4: Assigning the interface addresses}

\subsection{Static addresses on R2, R3, R4 and R5}

A static address is assigned to an interface with the command \texttt{ipv6 address}, followed by the address and the prefix length. Three more commands are used in every interface block:

\begin{itemize}
\item \texttt{interface FastEthernet0/0} enters the configuration of that interface. Everything that follows, until the next \texttt{interface} or the end of the block, applies to it.
\item \texttt{ipv6 address 2001:1000:0:12::2/64} assigns the global unicast address. The \texttt{/64} indicates that the first 64 bits identify the subnet and the last 64 bits identify the interface.
\item \texttt{no shutdown} activates the interface. Cisco interfaces start in the administratively down state, so this command is needed on every interface. An interface that is still administratively down appears as \texttt{administratively down/down} in the verification commands.
\end{itemize}

On R2 the three interfaces are configured as follows:

\begin{codeblock}
configure terminal
interface FastEthernet0/0
 ipv6 address 2001:1000:0:12::2/64
 no shutdown
interface FastEthernet0/1
 ipv6 address 2001:1000:0:23::2/64
 no shutdown
interface FastEthernet1/0
 ipv6 address 2001:1000:0:24::2/64
 no shutdown
end
\end{codeblock}

R3 has two interfaces, toward R2 and toward R5:

\begin{codeblock}
configure terminal
interface FastEthernet0/0
 ipv6 address 2001:1000:0:23::3/64
 no shutdown
interface FastEthernet0/1
 ipv6 address 2001:1000:0:35::3/64
 no shutdown
end
\end{codeblock}

R4 has two interfaces, toward R2 and toward R5:

\begin{codeblock}
configure terminal
interface FastEthernet0/0
 ipv6 address 2001:1000:0:24::4/64
 no shutdown
interface FastEthernet0/1
 ipv6 address 2001:1000:0:45::4/64
 no shutdown
end
\end{codeblock}

R5 has the reversed order described in Section~\ref{sec:lab01-topology}: FastEthernet0/1 toward R3 and FastEthernet0/0 toward R4.

\begin{codeblock}
configure terminal
interface FastEthernet0/1
 ipv6 address 2001:1000:0:35::5/64
 no shutdown
interface FastEthernet0/0
 ipv6 address 2001:1000:0:45::5/64
 no shutdown
end
\end{codeblock}

\subsection{Autoconfiguration on R1}
\label{sec:lab01-autoconfig}

R1 must not have a manually configured address. The handout requires that its address be built automatically through the Router Advertisement procedure. The only command needed is \texttt{ipv6 address autoconfig}:

\begin{codeblock}
configure terminal
interface FastEthernet0/0
 ipv6 address autoconfig
 no shutdown
end
\end{codeblock}

With this command R1 creates, on its interface, an address that combines the prefix announced by R2 with an interface identifier generated locally. R2 sends the prefix in the Router Advertisement messages that it transmits periodically on its interfaces, so R1 does not need to know the prefix in advance. The mechanism is explained in detail in Section~\ref{sec:lab01-rs-ra}.

The autoconfiguration depends on R2: R1 learns the prefix only if R2 already has an IPv6 address on the interface that faces R1. For this reason the verification of R1 is done after R2 has been configured.

\subsection{EUI-64 interface identifiers}
\label{sec:lab01-eui64}

The interface identifier of R1 is \texttt{C201:6FF:FE13:0} in the compressed notation of IOS. Its full 64-bit value is \texttt{C2:01:06:FF:FE:13:00:00}. The construction is the following:

\begin{enumerate}
\item Take the 48-bit MAC address of the interface, for example \texttt{C0:01:06:13:00:00}.
\item Insert the two bytes \texttt{FF:FE} between the third and the fourth byte, which gives \texttt{C0:01:06:FF:FE:13:00:00}.
\item Invert the universal/local bit of the first byte, which has value \texttt{0x02}. The byte \texttt{C0} becomes \texttt{C2}.
\end{enumerate}

The result is the interface identifier \texttt{C2:01:06:FF:FE:13:00:00}, which is written in IPv6 notation as \texttt{C201:06FF:FE13:0000} and abbreviated as \texttt{C201:6FF:FE13:0}. This is why the same MAC address always generates the same interface identifier.

\subsection{Switched interfaces and the \texttt{switchport} command}
\label{sec:lab01-switchport}

The routers of the laboratory have, besides the four FastEthernet0/x and 1/0 interfaces, a set of FastEthernet1/x interfaces (1/0 to 1/15) and a \texttt{Vlan1} interface. The \texttt{show ipv6 interface brief} output of R1 lists all of them. These are the ports of a switch module built into the router. Such ports work at layer 2 by default, which means that they forward Ethernet frames inside a VLAN and do not have IP addresses.

A layer 2 port is the normal behavior of a switch port. The command \texttt{switchport} explicitly configures a port as a layer 2 port. The opposite command, \texttt{no switchport}, turns a port into a routed port (layer 3), which is an interface that can have an IP address and route packets. Routed ports are the normal way to connect routers.

The handout explains that FastEthernet1/x interfaces are switched and must not be configured like regular router interfaces. Its example converts the interface to a routed port before assigning the address:

\begin{codeblock}
configure terminal
interface FastEthernet1/0
 no switchport
 ipv6 address 2001:1000:0:24::2/64
 no shutdown
exit
\end{codeblock}

The handout recommends applying \texttt{no switchport} before \texttt{ipv6 address} on these interfaces. In the run documented here, the link R2--R4 used R2 FastEthernet1/0 and the address was accepted and the interface came up without this command. The safe practice, however, is to follow the handout: if an interface shows \texttt{up/down} or the \texttt{ipv6 address} command returns an error, add \texttt{no switchport} and check again.

The topology contains no separate second switch. The only switching involved is the internal switch module of R2 and R1, and only one of its ports (R2 FastEthernet1/0) is used in the addressing plan.

\section{Step 5: Verifying the interfaces}

The command \texttt{show ipv6 interface brief} gives, for each interface, its state and its IPv6 addresses. It is run on every router after its configuration:

\begin{codeblock}
show ipv6 interface brief
\end{codeblock}

The first value in square brackets is the state of the physical and link layers, the second is the state of the protocol. \texttt{up/up} means that the interface is connected and working. \texttt{administratively down/down} means that \texttt{no shutdown} was not given. \texttt{up/down} on an interface that has no cable attached is normal, and it is the state of the unused FastEthernet1/x ports.

Each active interface shows at least one line starting with \texttt{FE80::}. This is the link-local address, which every IPv6 interface creates by itself and which is used on the link only, for control messages and for the next hop of many routing protocols. Below it, the global unicast address appears, if configured or autoconfigured.

The output of R1, after the autoconfiguration, is:

\begin{codeblock}
R1#show ipv6 interface brief
FastEthernet0/0          [up/up]
    FE80::C201:6FF:FE13:0
    2001:1000:0:12:C201:6FF:FE13:0
FastEthernet0/1          [administratively down/down]
FastEthernet1/0          [up/down]
...
\end{codeblock}

The address \texttt{2001:1000:0:12:C201:6FF:FE13:0} is the global address generated by R1. Its first 64 bits are the prefix \texttt{2001:1000:0:12::/64}, learned from R2; its last 64 bits, \texttt{C201:6FF:FE13:0}, are the interface identifier generated from the MAC address of the interface (Section~\ref{sec:lab01-eui64}). The output of R2 confirms the three static addresses configured in Step 4:

\begin{codeblock}
R2#show ipv6 interface brief
FastEthernet0/0          [up/up]
    FE80::C202:6FF:FE2A:0
    2001:1000:0:12::2
FastEthernet0/1          [up/up]
    FE80::C202:6FF:FE2A:1
    2001:1000:0:23::2
FastEthernet1/0          [up/up]
    FE80::C202:6FF:FE2A:F100
    2001:1000:0:24::2
\end{codeblock}

The same check on R3, R4 and R5 shows \texttt{up/up} on the two connected interfaces, with the addresses of Table~\ref{tab:lab01-plan}. If any address is missing or any interface is not \texttt{up/up}, the fix is to check the configuration of that interface and to repeat \texttt{no shutdown} if needed.

\section{Step 6: Static routing}

Static routing does not adapt to failures. If a link such as R2--R3 goes down, the routes that use it stop working until the administrator changes them by hand.

\subsection{How an IPv6 router chooses a route}

A router forwards each packet according to its routing table. For each packet, the router looks for the entry whose prefix matches the destination address with the largest number of bits. This rule is called \emph{longest prefix match}. The mechanism is the same as in IPv4, with one difference: IPv6 addresses have 128 bits, so the comparison covers more bits.

Each entry tells the router which interface to use and, for a remote network, which \emph{next hop} to send the packet to. The next hop is the address of the neighboring router that is closer to the destination. In a static route, the administrator writes both the destination prefix and the next hop.

The default route, written \texttt{::/0}, matches every address. It is used when no other entry matches, and it is the usual solution for a router that has only one way out, like R1.

\subsection{Routes must work in both directions}

A static route is needed in both directions. When R1 sends a packet to R5, R2 forwards it toward R5, but R5 must also send the reply to R1. If R5 has no route to the R1--R2 subnet, the reply is dropped and the \texttt{ping} fails, even if the packet reached its destination. This is the most frequent cause of failure in this laboratory, and it is why every router must know the path back to the subnet \texttt{2001:1000:0:12::/64}.

\subsection{The static routes of each router}

Table~\ref{tab:lab01-routes} lists the routes of the network. The reason column explains which branch of the topology each route uses.

\begin{table}[h]
\centering
\begin{tabular}{@{}lllp{5.8cm}@{}}
\toprule
\textbf{Router} & \textbf{Destination} & \textbf{Next hop} & \textbf{Reason} \\
\midrule
R1 & \texttt{::/0} & \texttt{2001:1000:0:12::2} & Everything goes to R2 \\
\midrule
R2 & \texttt{2001:1000:0:35::/64} & \texttt{2001:1000:0:23::3} & R5 reached through R3 \\
R2 & \texttt{2001:1000:0:45::/64} & \texttt{2001:1000:0:24::4} & R5 reached through R4 \\
\midrule
R3 & \texttt{2001:1000:0:12::/64} & \texttt{2001:1000:0:23::2} & Return path to R1 \\
R3 & \texttt{2001:1000:0:24::/64} & \texttt{2001:1000:0:23::2} & R4 reached through R2 \\
R3 & \texttt{2001:1000:0:45::/64} & \texttt{2001:1000:0:35::5} & R4--R5 subnet reached through R5 \\
\midrule
R4 & \texttt{2001:1000:0:12::/64} & \texttt{2001:1000:0:24::2} & Return path to R1 \\
R4 & \texttt{2001:1000:0:23::/64} & \texttt{2001:1000:0:24::2} & R3 reached through R2 \\
R4 & \texttt{2001:1000:0:35::/64} & \texttt{2001:1000:0:45::5} & R3--R5 subnet reached through R5 \\
\midrule
R5 & \texttt{2001:1000:0:12::/64} & \texttt{2001:1000:0:35::3} & Return path to R1 through R3 \\
R5 & \texttt{2001:1000:0:23::/64} & \texttt{2001:1000:0:35::3} & R2--R3 subnet reached through R3 \\
R5 & \texttt{2001:1000:0:24::/64} & \texttt{2001:1000:0:45::4} & R2--R4 subnet reached through R4 \\
\bottomrule
\end{tabular}
\caption{Static routes of the laboratory.}
\label{tab:lab01-routes}
\end{table}

Some destinations are directly connected to a router and need no static route. For example, R2 is directly connected to \texttt{2001:1000:0:12::/64}, \texttt{2001:1000:0:23::/64} and \texttt{2001:1000:0:24::/64}, and it does not need static entries for them.

\subsection{Entering the routes}

The command has the form \texttt{ipv6 route <destination/length> <next hop>}, typed in global configuration mode. Here are the routes of R1 and R2:

\begin{codeblock}
configure terminal
ipv6 route ::/0 2001:1000:0:12::2
end
\end{codeblock}

\begin{codeblock}
configure terminal
ipv6 route 2001:1000:0:35::/64 2001:1000:0:23::3
ipv6 route 2001:1000:0:45::/64 2001:1000:0:24::4
end
\end{codeblock}

R4 and R5 use the same form:

\begin{codeblock}
configure terminal
ipv6 route 2001:1000:0:12::/64 2001:1000:0:24::2
ipv6 route 2001:1000:0:23::/64 2001:1000:0:24::2
ipv6 route 2001:1000:0:35::/64 2001:1000:0:45::5
end
\end{codeblock}

\begin{codeblock}
configure terminal
ipv6 route 2001:1000:0:12::/64 2001:1000:0:35::3
ipv6 route 2001:1000:0:23::/64 2001:1000:0:35::3
ipv6 route 2001:1000:0:24::/64 2001:1000:0:45::4
end
\end{codeblock}

The routes of R3 were entered in the same form, but the form that finally worked is the one with the outgoing interface, explained in the next subsection:

\begin{codeblock}
configure terminal
ipv6 route 2001:1000:0:12::/64 FastEthernet0/0 2001:1000:0:23::2
ipv6 route 2001:1000:0:24::/64 FastEthernet0/0 2001:1000:0:23::2
ipv6 route 2001:1000:0:45::/64 FastEthernet0/1 2001:1000:0:35::5
end
\end{codeblock}

The commands \texttt{show running-config | include ipv6 route} and \texttt{show ipv6 route static} check the routes in two different ways. The first shows the routes written in the configuration; the second shows the routes that are actually installed in the routing table, which are the ones that forward packets.

\subsection{A route that was configured but not installed}
\label{sec:lab01-static-issue}

On R3, the three routes entered with the next hop only appeared in \texttt{show running-config}, but \texttt{show ipv6 route static} showed no static entries. The table contained only the connected and local entries, and the ping from R1 failed. A route that is in the configuration but not in the table does not forward any packet.

The fix was to remove the three routes and to enter them again with the outgoing interface:

\begin{codeblock}
configure terminal
no ipv6 route 2001:1000:0:12::/64 2001:1000:0:23::2
ipv6 route 2001:1000:0:12::/64 FastEthernet0/0 2001:1000:0:23::2
no ipv6 route 2001:1000:0:24::/64 2001:1000:0:23::2
ipv6 route 2001:1000:0:24::/64 FastEthernet0/0 2001:1000:0:23::2
no ipv6 route 2001:1000:0:45::/64 2001:1000:0:35::5
ipv6 route 2001:1000:0:45::/64 FastEthernet0/1 2001:1000:0:35::5
end
\end{codeblock}

The \texttt{no} form removes a route and must repeat the same destination and next hop of the original command. After the change, \texttt{show ipv6 route static} on R3 listed the three routes with their outgoing interfaces:

\begin{codeblock}
R3#show ipv6 route static
S   2001:1000:0:12::/64 [1/0]
     via 2001:1000:0:23::2, FastEthernet0/0
S   2001:1000:0:24::/64 [1/0]
     via 2001:1000:0:23::2, FastEthernet0/0
S   2001:1000:0:45::/64 [1/0]
     via 2001:1000:0:35::5, FastEthernet0/1
\end{codeblock}

The handout does not explain this behavior, and the exact cause was not identified during the laboratory. The practical rule is to check the routing table with \texttt{show ipv6 route static} after every route entered, and to use the form with the outgoing interface whenever a route does not appear.

\section{Step 7: Verifying the routing table}

The routing table is shown with \texttt{show ipv6 route}. On R1, after the configuration, the output is:

\begin{codeblock}
R1#show ipv6 route
S   ::/0 [1/0]
     via 2001:1000:0:12::2
C   2001:1000:0:12::/64 [0/0]
     via ::, FastEthernet0/0
L   2001:1000:0:12:C201:6FF:FE13:0/128 [0/0]
     via ::, FastEthernet0/0
L   FE80::/10 [0/0]
     via ::, Null0
L   FF00::/8 [0/0]
     via ::, Null0
\end{codeblock}

The codes at the beginning of each entry have a precise meaning:

\begin{itemize}
\item \textbf{C} (connected): the subnet is directly attached to an interface of the router.
\item \textbf{L} (local): the address of the router itself on that interface, with prefix length \texttt{/128}. It is not a route to another network.
\item \textbf{S} (static): the route was configured by the administrator.
\item The \texttt{FE80::/10} and \texttt{FF00::/8} entries pointing to \texttt{Null0} are built-in system entries. They discard traffic for link-local and multicast destinations that the router does not handle, and they do not need to be configured.
\end{itemize}

The number in brackets, for example \texttt{[1/0]}, is the administrative distance followed by the metric. The administrative distance expresses how much a router trusts a source of routes; static routes have the value 1, which is lower than most dynamic routing protocols. The metric is used to choose between several routes of the same source. In this laboratory each destination has only one static route, so the values have no effect on the choice.

\section{Step 8: Testing connectivity with ping}

The final test is a \texttt{ping} from R1 to one of the addresses of R5:

\begin{codeblock}
R1#ping 2001:1000:0:45::5
\end{codeblock}

The \texttt{ping} command sends ICMPv6 echo requests (message type 128) to the destination, and it waits for the echo replies (message type 129). Each request that receives a reply is printed as an exclamation mark, and each request without a reply is printed as a dot after the timeout.

The first attempt, made when only the default route of R1 was configured, failed:

\begin{codeblock}
R1#ping 2001:1000:0:45::5
Type escape sequence to abort.
Sending 5, 100-byte ICMP Echos to 2001:1000:0:45::5, timeout is 2 seconds:
.....
Success rate is 0 percent (0/5)
\end{codeblock}

The default route of R1 was correct: the request left R1 and reached R2. The failure came from the routers beyond R2, which at that point had few or no static routes. After the routes of R2, R4 and R5 were entered, the ping still failed, because the routes of R3 were not installed in its routing table (Section~\ref{sec:lab01-static-issue}). Once those routes were installed with the outgoing interface, the ping succeeded:

\begin{codeblock}
R1#ping 2001:1000:0:45::5
Type escape sequence to abort.
Sending 5, 100-byte ICMP Echos to 2001:1000:0:45::5, timeout is 2 seconds:
!!!!!
Success rate is 100 percent (5/5)
\end{codeblock}

A successful ping proves that the routing works in both directions along one path, and that every router of the path has the required entries. The handout suggests testing one of the two addresses of R5. Testing also \texttt{2001:1000:0:35::5} is a useful additional check, because that address is reached through the other branch of the network.

\section{Step 9: Traffic analysis with Wireshark}

The last part of the laboratory observes, in a packet capture, the control messages of the protocol. The capture is started in GNS3 by right-clicking the link to observe and choosing the capture option; GNS3 opens the traffic in Wireshark. The display filters used here select the messages by ICMPv6 type.

The two questions of the handout ask for a description of the principles of operation of two procedures. The answers below explain the procedures and the messages to look for in the capture.

\subsection{Autoconfiguration with RS and RA (review question 4)}
\label{sec:lab01-rs-ra}

The autoconfiguration of R1 is based on two ICMPv6 messages, Router Solicitation and Router Advertisement. The sequence is the following.

\begin{enumerate}
\item \textbf{Creation of the link-local address.} As soon as R1 enables the interface, it creates an address in the \texttt{FE80::/64} range from its MAC address. This address is needed to send control messages on the link, and it exists even before any prefix is known.
\item \textbf{Router Solicitation (RS, type 133).} R1 sends an RS to the all-routers multicast address \texttt{ff02::2}. The message asks any router on the link to announce the prefixes available. The source of the RS is the link-local address of R1, and the destination MAC address is the multicast address \texttt{33:33:00:00:00:02}.
\item \textbf{Router Advertisement (RA, type 134).} R2 answers with an RA sent to the all-nodes multicast address \texttt{ff02::1}. The RA contains the prefix \texttt{2001:1000:0:12::/64} in a \emph{Prefix Information} option. Two flags of this option matter for autoconfiguration: the \textbf{L} flag says that the prefix is on-link, and the \textbf{A} flag says that hosts may use the prefix to build their addresses. R2 also sends RA messages periodically, without waiting for an RS, so the prefix is learned even when R1 starts with the interface already up.
\item \textbf{Construction of the address.} R1 takes the 64 bits of the prefix and appends the 64-bit interface identifier, obtained from its MAC address with the EUI-64 method. The result is \texttt{2001:1000:0:12:C201:6FF:FE13:0/64}.
\item \textbf{Duplicate Address Detection (DAD).} Before using the new address, R1 checks that no other node on the link has it. It sends a Neighbor Solicitation (Section~\ref{sec:lab01-ns-na}) to the solicited-node address of the new address, with the unspecified source address \texttt{::}. If a reply arrives, another node is using the address and R1 does not assign it.
\end{enumerate}

In the capture, the RS and the RA are the two messages with ICMPv6 type 133 and 134, and the filter \texttt{icmpv6.type == 133 or icmpv6.type == 134} shows both. Expanding the RA shows the Prefix Information option with the prefix and the flags described above.

The interface identifier is built with the \emph{EUI-64} method, described in Section~\ref{sec:lab01-eui64}. Because the method is reversible, the MAC address of the interface can be recovered from the interface identifier of R1: the value \texttt{C201:06FF:FE13:0000} gives the MAC address \texttt{C0:01:06:13:00:00}, which can be checked with \texttt{show interfaces FastEthernet0/0}.

\subsection{Neighbor discovery and MAC address discovery (review question 5)}
\label{sec:lab01-ns-na}

In IPv4 the MAC address of a neighbor is discovered with ARP, which sends a broadcast request to all hosts. In IPv6 the same job is done by Neighbor Discovery, which uses ICMPv6 messages and multicast addresses. When R1 pings R5, the packet must travel first to R2, the next hop of R1. R1 therefore needs the MAC address of R2, and it discovers it as follows:

\begin{enumerate}
\item R1 sends a \textbf{Neighbor Solicitation (NS, type 135)} asking which MAC address corresponds to the IPv6 address \texttt{2001:1000:0:12::2} of R2.
\item The NS is not sent to all nodes. It is sent to the \textbf{solicited-node multicast address} of the target, \texttt{ff02::1:ff00:2}. This address is built from the prefix \texttt{ff02::1:ff00:0/104} and the last 24 bits of the target address (\texttt{00:00:02}). Only the nodes whose addresses end with the same 24 bits join this group, so in practice only R2 receives the message.
\item The Ethernet frame carrying the NS uses the destination MAC address \texttt{33:33:ff:00:00:02}. The prefix \texttt{33:33} identifies IPv6 multicast frames, and the last 32 bits are copied from the IPv6 multicast address. Network cards filter these frames, so the other hosts of the link do not process them.
\item R2 answers with a \textbf{Neighbor Advertisement (NA, type 136)}, sent in unicast to R1. The NA contains the MAC address of R2 in the \emph{Target Link-Layer Address} option.
\item R1 stores the pair IPv6 address and MAC address in its \textbf{neighbor cache}. The cache is shown with \texttt{show ipv6 neighbors}, and it plays the role of the ARP table of IPv4.
\end{enumerate}

The same process happens for every hop of the route, in both directions. This is why the capture of the ping shows NS and NA messages not only between R1 and R2, but also between the other routers of the path. The NA flags describe the role of the sender: the \textbf{R} flag is set when the sender is a router, the \textbf{S} flag when the advertisement answers a solicitation, and the \textbf{O} flag when the advertisement may overwrite an entry already in the cache.

The table below summarizes the differences between the two protocols.

\begin{table}[h]
\centering
\begin{tabular}{@{}lp{5.6cm}p{5.6cm}@{}}
\toprule
\textbf{Aspect} & \textbf{IPv4 (ARP and DHCP)} & \textbf{IPv6 (Neighbor Discovery)} \\
\midrule
Address resolution & ARP request, sent to the broadcast address & Neighbor Solicitation, sent to a solicited-node multicast address \\
Address resolution reply & ARP reply & Neighbor Advertisement \\
Router discovery & Default gateway configured manually or through DHCP & Router Solicitation and Router Advertisement \\
Protocol & Separate from IP & Part of ICMPv6 \\
\bottomrule
\end{tabular}
\caption{Comparison between ARP in IPv4 and Neighbor Discovery in IPv6.}
\end{table}

\section{Answers to the review questions}

The handout asks five questions. They are answered here with the information obtained in the laboratory, so that the answers can be checked against the configuration.

\textbf{Question 1.} The prefixes and the interface addresses are those of Table~\ref{tab:lab01-plan}.

\textbf{Question 2.} The static routes of each router are those of Table~\ref{tab:lab01-routes}.

\textbf{Question 3.} The global unicast address autoconfigured on R1 FastEthernet0/0 is \texttt{2001:1000:0:12:C201:6FF:FE13:0}.

\textbf{Question 4.} The autoconfiguration works as described in Section~\ref{sec:lab01-rs-ra}: R1 sends a Router Solicitation to the all-routers multicast address, R2 answers with a Router Advertisement to the all-nodes multicast address, the advertisement carries the prefix \texttt{2001:1000:0:12::/64}, and R1 appends its 64-bit interface identifier to it. Before using the address, R1 performs Duplicate Address Detection.

\textbf{Question 5.} When R1 pings R5, it needs the MAC address of its next hop. It sends a Neighbor Solicitation to the solicited-node multicast address of the next hop, receives a Neighbor Advertisement with the MAC address of the next hop, and stores the result in its neighbor cache, as described in Section~\ref{sec:lab01-ns-na}.

The answers to questions 4 and 5 must be supported by the Wireshark captures of the RS, RA, NS and NA messages, which are required by the handout.

\section{Troubleshooting}

Table~\ref{tab:lab01-trouble} collects the problems met during the laboratory, with their causes and solutions.

\begin{table}[h]
\centering
\begin{tabular}{@{}p{4.3cm}p{5.4cm}p{5.6cm}@{}}
\toprule
\textbf{Symptom} & \textbf{Cause} & \textbf{Solution} \\
\midrule
\texttt{\% Incomplete command} after \texttt{ipv6 route} & The address was typed with a single colon where two were needed, for example \texttt{2001:1000:0:12:2} instead of \texttt{2001:1000:0:12::2} & Type the address with the double colon that represents the zero hextets \\
\midrule
\texttt{\% Invalid input detected at '\^{}' marker} & A word or a character is not valid in the current mode, or a command was typed in the wrong mode & Read the prompt, check the mode, and type the command again \\
\midrule
A route is in \texttt{show running-config} but not in \texttt{show ipv6 route static} & The route was not installed in the routing table & Enter the route again with the outgoing interface, as in Step 6 \\
\midrule
Ping fails with \texttt{Success rate is 0 percent} & A router on the path has no route toward the destination or toward the source of the request & Check \texttt{show ipv6 route} on every router of the path, not only on R1 \\
\midrule
An interface stays \texttt{administratively down/down} & The command \texttt{no shutdown} was not given & Enter the interface and give \texttt{no shutdown} \\
\midrule
R1 has no global address & The Router Advertisement has not been received yet, or R2 has no address on the interface facing R1 & Wait a few seconds and repeat \texttt{show ipv6 interface brief}; check the address of R2 Fa0/0 \\
\midrule
An interface on FastEthernet1/x does not accept the address & The port is a switched port & Apply \texttt{no switchport} before \texttt{ipv6 address}, as in Section~\ref{sec:lab01-switchport} \\
\bottomrule
\end{tabular}
\caption{Problems met during the laboratory.}
\label{tab:lab01-trouble}
\end{table}

\section{Glossary of the new concepts}
\label{sec:lab01-glossary}

Table~\ref{tab:lab01-glossary} collects the main terms introduced in this laboratory.

\begin{table}[h]
\centering
\begin{tabular}{@{}p{4.2cm}p{11cm}@{}}
\toprule
\textbf{Term} & \textbf{Meaning} \\
\midrule
Hextet & One of the eight 16-bit groups of an IPv6 address, written as four hexadecimal digits. \\
Prefix length & The number after the slash, which says how many leading bits identify the network (for example \texttt{/56} or \texttt{/64}). \\
Subnet & The set of addresses that share the same prefix of length \texttt{/64} on a link. \\
Interface identifier & The last 64 bits of an address, which identify the interface inside its subnet. \\
Link-local address & An address in \texttt{FE80::/10}, created automatically on each interface and valid only on the local link. \\
Global unicast address & An address routable on the whole network, such as \texttt{2001:1000:0:12::2}. \\
EUI-64 & The method that builds a 64-bit interface identifier from a 48-bit MAC address: the bytes \texttt{FF:FE} are inserted in the middle, and the universal/local bit of the first byte (value \texttt{0x02}) is inverted. \\
SLAAC & Stateless Address Autoconfiguration. The host builds its address from the prefix in the Router Advertisements and from its own interface identifier, without a server that keeps track of the assignments. \\
Router Solicitation (RS) & ICMPv6 type 133. A host asks the routers on the link to send a Router Advertisement. \\
Router Advertisement (RA) & ICMPv6 type 134. A router announces the prefixes and the configuration parameters of the link. \\
Duplicate Address Detection (DAD) & The check that a new address is not already used by another node, performed with a Neighbor Solicitation before the address is used. \\
Neighbor Solicitation (NS) & ICMPv6 type 135. A node asks for the link-layer address of a neighbor, or checks that an address is free. \\
Neighbor Advertisement (NA) & ICMPv6 type 136. A node answers a solicitation or announces its link-layer address. \\
Solicited-node multicast address & A multicast address built from the last 24 bits of a unicast address, used to reach the nodes that may own that address without sending to everyone. \\
Neighbor cache & The table of IPv6 addresses and MAC addresses learned through Neighbor Discovery. It is the IPv6 equivalent of the ARP table. \\
Next hop & The neighboring router to which a packet is sent in order to reach its destination. \\
Default route & The route \texttt{::/0}, used when no other route matches the destination. \\
Longest prefix match & The rule by which a router chooses the routing entry whose prefix has the largest number of matching bits with the destination. \\
Connected, local and static routes & Codes \texttt{C}, \texttt{L} and \texttt{S} in the routing table: a directly attached subnet, the router's own address, and a manually configured route. \\
Routed port (\texttt{no switchport}) & A layer 3 interface that can have an IP address and route packets. \\
Switched port (\texttt{switchport}) & A layer 2 interface that forwards Ethernet frames inside a VLAN and has no IP address. \\
Administratively down & The state of an interface that has been disabled in the configuration, until \texttt{no shutdown} is given. \\
\bottomrule
\end{tabular}
\caption{Glossary of the new concepts.}
\label{tab:lab01-glossary}
\end{table}

\section{Command reference}
\label{sec:lab01-commands}

Table~\ref{tab:lab01-commands} collects the commands used in the laboratory, grouped by purpose.

\begin{table}[h]
\centering
\begin{tabular}{@{}p{6.6cm}p{8.6cm}@{}}
\toprule
\textbf{Command} & \textbf{Effect} \\
\midrule
\texttt{enable} & Enters privileged mode (prompt \texttt{R1\#}). \\
\texttt{configure terminal} & Enters global configuration mode (prompt \texttt{R1(config)\#}). \\
\texttt{end} & Returns to privileged mode. \\
\texttt{reload} & Restarts the router. Unsaved changes are discarded if the save prompt is answered \texttt{no}. \\
\texttt{erase startup-config} & Deletes the saved configuration. Used with \texttt{reload} for a full reset. \\
\texttt{ipv6 unicast-routing} & Enables IPv6 and IPv6 packet forwarding on the router. \\
\texttt{interface <name>} & Enters the configuration of an interface. \\
\texttt{ipv6 address <address/length>} & Assigns a static global address to the interface. \\
\texttt{ipv6 address autoconfig} & Enables stateless autoconfiguration on the interface. \\
\texttt{no shutdown} & Activates the interface. \\
\texttt{no switchport} & Converts a switched port into a routed port. \\
\texttt{switchport} & Configures a port as a layer 2 (switched) port. \\
\texttt{ipv6 route <prefix/length> <next hop>} & Adds a static route with the given next hop. \\
\texttt{ipv6 route <prefix/length> <interface> <next hop>} & Adds a static route with the given next hop and outgoing interface. \\
\texttt{ipv6 route ::/0 <next hop>} & Adds the default route. \\
\texttt{no ipv6 route <prefix/length> <next hop>} & Removes a static route, using the same destination and next hop. \\
\texttt{show ipv6 interface brief} & Shows the state and the addresses of all interfaces. \\
\texttt{show ipv6 route} & Shows the whole routing table. \\
\texttt{show ipv6 route static} & Shows only the static routes that are installed. \\
\texttt{show running-config | include ipv6} & Shows the lines containing \texttt{ipv6} in the current configuration. \\
\texttt{show ipv6 neighbors} & Shows the neighbor cache (IPv6 and MAC addresses). \\
\texttt{ping <address>} & Sends ICMPv6 echo requests to the address and reports the success rate. \\
\bottomrule
\end{tabular}
\caption{Commands used in the laboratory.}
\label{tab:lab01-commands}
\end{table}

